\chapter{Week 3: Types and a First Interpreter}
\label{ch:week3}

Lecture 3: September 22, 2026. Recitation 3: September 25, 2026.

Supplementary reading:
\begin{itemize}
  \item \cite{PLAI} Chapters 2 (Everything (We Will Say) About Parsing) and 3 (A First Look at Interpretation).
  \item \cite{PLAI} Chapter 15 (Checking Program Invariants Statically: Types), first sections, and the \code{plai-typed} documentation.
  \item \cite{TAPL} Chapter 1 (what type systems are for), optional.
\end{itemize}

From this week on, code is written in \code{\#lang plai-typed}, the language of the textbook: Racket with a type checker that reads the whole program before anything runs. Lecture~3 introduces its types, defines new types with \code{define-type}, and uses them to write the first interpreter of the course, a calculator. Recitation~3 redoes Recitation~2 with types and then studies Racket's \code{match} in more depth. This chapter covers:
\begin{itemize}
  \item \Cref{sec:types}: types, type inference and annotations, function types, homogeneous lists, pairs, and polymorphism;
  \item \Cref{sec:definetype}: \code{define-type}, \code{type-case}, and exhaustiveness checking;
  \item \Cref{sec:calc}: quote and s-expressions, concrete vs.\ abstract syntax, and the calculator's parser and interpreter;
  \item \Cref{sec:rec3}: Recitation~3, typed Boolean expressions and pattern matching with \code{match}.
\end{itemize}

\section{Types in plai-typed}
\label{sec:types}

\code{plai-typed} is installed with \code{raco pkg install plai-typed}, and a file starts with \code{\#lang plai-typed}. It looks like Racket, but a type checker reads the whole program first, and a program that fails the check does not run at all.

\subsection{Every Value Has a Type}

At the REPL, every answer comes with two lines: the type the checker worked out, starting with \code{-}, and then the value.
\begin{lstlisting}[style=repl]
> 3
- number
3
> (+ 3 4)
- number
7
> "hello"
- string
"hello"
> (+ 3 "hello")
typecheck failed: number vs. string
\end{lstlisting}
This is the idea of the lecture in miniature: the type is known \emph{before} the value is computed. The checker does not run \rkt|(+ 3 4)| to find out it is a number; it knows because \code{+} takes numbers and returns a number.

\subsection{Type Inference and Annotations}

No types need to be written. Here the checker sees \rkt|(+ 3 x)|, concludes that \code{x} must be a number, and that the result is one too:
\begin{lstlisting}
(define (add3 x)
  (+ 3 x))
\end{lstlisting}

\begin{definition}[Type inference]
\emph{Type inference} is the checker working out the types that the programmer did not write down.
\end{definition}

Types can also be written, as many or as few as we like. An argument becomes \rkt|(x : number)|, and the return type goes after the argument list. Whatever is left out is inferred:
\begin{lstlisting}
(define (add3-v2 (x : number))            ; return type inferred
  (+ 3 x))

(define (add3-v3 (x : number)) : number   ; both annotated
  (+ 3 x))
\end{lstlisting}

\begin{pitfall}[Spaces around the colon]
\code{(x: number)} is read as one identifier named \code{x:}, followed by \code{number}, and is rejected with \code{expected either <id> or `[<id> : <type>]'}. Write \code{(x : number)}.
\end{pitfall}

An annotation is a claim the checker verifies, not a hint it trusts. This one is wrong, since \rkt|(+ 3 x)| is a number, so the definition is rejected:
\begin{lstlisting}
(define (add3-v4 (x : number)) : string
  (+ 3 x))
;; typecheck failed: number vs. string
\end{lstlisting}
Either way the program is checked before it runs: a call like \rkt|(add3-v3 "hello")| never gets the chance to fail at run time, because the whole file is rejected.

If the checker can infer the types, why write them? Two reasons from lecture:
\begin{itemize}
  \item \textbf{Documentation.} The reader sees what the function takes and returns without reading its body.
  \item \textbf{Mental model.} Writing the type first makes you decide what the function is supposed to do before you write how it does it.
\end{itemize}

\subsection{Function Types}
\label{sec:fntypes}

Functions are values, so they have types. Read the arrow as ``takes \ldots{} and returns \ldots'':
\begin{lstlisting}[style=repl]
> add3
- (number -> number)
\end{lstlisting}
With two arguments, both go to the left of the arrow. Annotations can be mixed: in \code{myadd-v2}, \code{y} is unannotated but added to a number, so it is inferred to be one, and the two functions have the same type.
\begin{lstlisting}
(define (myadd x y)
  (+ x y))
(define (myadd-v2 (x : number) y) : number
  (+ x y))
\end{lstlisting}
\begin{lstlisting}[style=repl]
> myadd
- (number number -> number)
> myadd-v2
- (number number -> number)
\end{lstlisting}

A function that returns a function, \code{make-adder} from \Cref{sec:hof}, now has a type that says so:
\begin{lstlisting}
(define (myadd-v3 (x : number))
  (lambda (y) (+ x y)))
\end{lstlisting}
\begin{lstlisting}[style=repl]
> myadd-v3
- (number -> (number -> number))
> (myadd-v3 2)
- (number -> number)
> ((myadd-v3 2) 5)
- number
7
\end{lstlisting}
Compare the two types. \code{myadd} takes two numbers at once. \code{myadd-v3} takes \emph{one} number and returns a function that waits for the second, so \rkt|(myadd-v3 2)| is not a number yet but a function.

\subsection{Lists Are Homogeneous}

A list type says what its elements are, and every element must have that type:
\begin{lstlisting}[style=repl]
> (list 1 2 3 4)
- (listof number)
'(1 2 3 4)
> (cons 3 (cons 4 empty))
- (listof number)
'(3 4)
> (list 1 "hello")
typecheck failed: number vs. string
> (cons 3 4)
typecheck failed: (listof number) vs. number
\end{lstlisting}
Plain Racket accepts both of the last two. In \code{plai-typed} the second argument of \code{cons} must be a list, so bare pairs (\Cref{sec:lists}) are impossible. In exchange, list operations have precise types: the checker knows \rkt|(first (list 1 2 3 4))| is a \code{number} and \rkt|(rest (list 1 2 3 4))| a \code{(listof number)}, and likewise for \code{second} and \code{third}, without running anything.

\subsection{Pairs}

When two values of different types belong together, use a \emph{pair} (a tuple). Its type is written with \code{*}:
\begin{lstlisting}[style=repl]
> (pair 1 "hello")
- (number * string)
'#(1 "hello")
> (fst (pair 1 2))
- number
1
> (snd (pair 1 2))
- number
2
> (car (pair 1 2))
car: unbound identifier
\end{lstlisting}
A list has any length but one element type; a pair has exactly two slots, each with its own type. A pair is not a cons cell, and \code{plai-typed} provides no \code{car} or \code{cdr} at all: use \code{fst}/\code{snd} for pairs and \code{first}/\code{rest} for lists.

\subsection{Polymorphic Types}

Asking for the type of \code{first} shows something new:
\begin{lstlisting}[style=repl]
> first
- ((listof 'a) -> 'a)
> rest
- ((listof 'a) -> (listof 'a))
> cons
- ('a (listof 'a) -> (listof 'a))
> fst
- (('a * 'b) -> 'a)
> map
- (('a -> 'b) (listof 'a) -> (listof 'b))
\end{lstlisting}

\begin{definition}[Type variable, parametric polymorphism]
\code{'a} (read ``alpha'') is a \emph{type variable}: any type may stand in its place, but every \code{'a} in one type must stand for the \emph{same} type. A single definition that works uniformly at every type in this way is \emph{parametrically polymorphic}; other languages call it generics.
\end{definition}

So \code{first} takes a \code{(listof number)} to a \code{number} and a \code{(listof string)} to a \code{string}, but can never take a list of numbers to a string. Different letters allow different types: \code{fst} takes a pair of an \code{'a} and a \code{'b} and returns the \code{'a}. And the type of \code{map} alone says most of what \code{map} does: given a function from \code{'a} to \code{'b} and a list of \code{'a}, it produces a list of \code{'b}.

\subsection{Why Insist on Homogeneous Lists?}

Suppose \rkt|(list 1 "hello")| were allowed; call it \code{l}. What is the type of \rkt|(first l)|? It depends on which element happens to come first, and the checker cannot know that before the program runs. Then
\begin{itemize}
  \item every use of \code{first} would need a check at run time, which is exactly what a type checker is supposed to spare us, and
  \item \rkt|(map f l)| is only safe if \code{f} accepts \emph{both} numbers and strings, which \code{('a -> 'b)} cannot express.
\end{itemize}
One could design a type that lists every position, something like \code{(listof number string number)} for \rkt|(list 1 "hello" 2)|, but then every list function (\code{map}, \code{filter}, the folds) would need a far more complicated type. The simpler fix keeps lists homogeneous and instead makes \emph{one} type that covers all the kinds of things to be mixed. That is what \code{define-type} is for.

\section{\texttt{define-type} and \texttt{type-case}}
\label{sec:definetype}

\subsection{Defining a Type with Variants}

A shape is either a circle or a rectangle:
\begin{lstlisting}
(define-type Shape
  [circle (radius : number)]                       ; variant 1
  [rectangle (width : number) (height : number)])  ; variant 2
\end{lstlisting}
This defines one new type, \tyname{Shape}, with two \emph{variants}. Each variant has a constructor and typed fields:
\begin{lstlisting}[style=repl]
> (circle 3)
- Shape
(circle 3)
\end{lstlisting}
Now \rkt|(list (circle 1) (rectangle 2 3))| is fine: the elements are different kinds of shapes, but they have one type, so the list is a \code{(listof Shape)}.

Like \code{struct}, \code{define-type} also generates
\begin{itemize}
  \item predicates \code{circle?} and \code{rectangle?}, of type \code{(Shape -> boolean)}, and
  \item accessors \code{circle-radius}, \code{rectangle-width}, and \code{rectangle-height}, of type \code{(Shape -> number)}.
\end{itemize}
The predicates accept only shapes. In plain Racket \rkt|(circle? 5)| would return \code{\#f}; here it never runs:
\begin{lstlisting}[style=repl]
> (circle? 5)
typecheck failed: Shape vs. number
\end{lstlisting}

\subsubsection{A recursive type} A type may mention itself. The binary tree of \Cref{sec:structs}, now with types:
\begin{lstlisting}
(define-type Tree
  [emp]
  [node (n : number) (lt : Tree) (rt : Tree)])
\end{lstlisting}
A variant with no fields is still a constructor, so the empty tree is \rkt|(emp)|, with parentheses. \rkt|(node 3 (emp) (emp))| has type \tyname{Tree}.

\subsection{\texttt{struct} vs.\ \texttt{define-type}}

Both make compound values, but they promise different things:
\begin{center}
\begin{tabular}{@{}p{0.47\linewidth}p{0.47\linewidth}@{}}
\toprule
\code{struct} (\code{\#lang racket}) & \code{define-type} (\code{plai-typed}) \\
\midrule
one kind of value: a \code{bt} is always a \code{bt} & one \emph{type} with several variants: \code{circle} or \code{rectangle} \\
fields unchecked: \rkt|(bt 1 2 3)| is accepted, and fails later & every field typed: \rkt|(node 3 5 (emp))| is rejected before the program runs, at the line that is wrong \\
different structs cannot share a homogeneous list type & all variants fit in one \code{(listof Shape)} \\
\bottomrule
\end{tabular}
\end{center}

\subsubsection{Why only one \code{circle}?} No other type in the program may also have a variant named \code{circle}. The immediate reason is that \code{circle}, \code{circle?}, and \code{circle-radius} are ordinary definitions, so a second \code{define-type} would define them again (\code{identifier already defined}). The deeper reason is inference: when the checker sees \rkt|(circle 3)|, the constructor's name is all it has to go on. If two types both had a \code{circle} variant, \rkt|(circle 3)| could belong to either, and inference would be stuck.

\subsection{\texttt{type-case}: Taking a Value Apart}

A \code{define-type} value is built with a constructor; \code{type-case} takes it apart again. There is one clause per variant, and the names in parentheses are bound to that variant's fields, in order (the names need not match those in \code{define-type}; only the order matters):
\begin{lstlisting}
(define (area s)
  (type-case Shape s
    [circle (r) (* 3.14159 (* r r))]
    [rectangle (w h) (* w h)]))
\end{lstlisting}
\begin{lstlisting}[style=repl]
> (area (circle 3))
- number
28.27431
> area
- (Shape -> number)
\end{lstlisting}
The same function can be written with the generated predicates and accessors:
\begin{lstlisting}
(define (area-v2 s)
  (cond
    [(circle? s) (* 3.14159 (* (circle-radius s) (circle-radius s)))]
    [(rectangle? s) (* (rectangle-width s) (rectangle-height s))]))
\end{lstlisting}
It works, but \code{type-case} is shorter, names the fields, and lines the code up with the type: one clause per variant, just as list functions had one \code{cond} branch per kind of list.

\subsection{Exhaustiveness}

The real difference shows when a case is forgotten. \code{type-case} knows every variant of \tyname{Shape}, so leaving one out is rejected before the program runs:
\begin{lstlisting}
(define (area-v3 s)
  (type-case Shape s
    [circle (r) (* 3.14159 (* r r))]))
;; type-case: syntax error; probable cause: you did not include a case
;; for the rectangle variant, or no else-branch was present
\end{lstlisting}
\code{cond} knows nothing about \tyname{Shape}; it tries its tests in order. The same mistake with \code{cond} is accepted, works on circles, and fails only at run time when it meets a rectangle:
\begin{lstlisting}
(define (area-v4 s)
  (cond
    [(circle? s) (* 3.14159 (* (circle-radius s) (circle-radius s)))]))
\end{lstlisting}
\begin{lstlisting}[style=repl]
> (area-v4 (circle 1))
- number
3.14159
> (area-v4 (rectangle 2 3))
cond: no matching clause
\end{lstlisting}

\begin{keypoint}[What \code{type-case} buys]
\code{type-case} checks before the program runs that every variant is handled, and it gives each clause the variant's fields with their types. A missing case is a static error, not a bug waiting for the right input.
\end{keypoint}

\section{S-Expressions and the Calculator}
\label{sec:calc}

\subsection{Quote and Symbols}

The apostrophe \code{'}, read ``quote'', means ``do not evaluate this; keep it as data''. Without it, \code{x} is a variable to look up. With it, \code{'x} is the name itself, a \emph{symbol}:
\begin{lstlisting}[style=repl]
> x
x: free variable while typechecking in: x
> 'x
- symbol
'x
\end{lstlisting}
Quoting a parenthesized expression keeps the whole thing as data, without doing the addition:
\begin{lstlisting}[style=repl]
> '(1 (+ 2 3) 4)
- s-expression
'(1 (+ 2 3) 4)
\end{lstlisting}

\subsection{S-Expressions}

Look at what \code{'(+ 2 3)} contains: a symbol and two numbers. That is exactly the kind of mixed list that \Cref{sec:types} ruled out. So \code{plai-typed} gives all quoted data one type, \tyname{s-expression}, and provides functions to ask what a piece is and to convert it back out:
\begin{lstlisting}[style=repl]
> (s-exp->list '(+ 2 3))
- (listof s-expression)
'(+ 2 3)
> (first (s-exp->list '(+ 2 3)))
- s-expression
'+
> (s-exp->symbol (first (s-exp->list '(+ 2 3))))
- symbol
'+
> (s-exp->number (second (s-exp->list '(+ 2 3))))
- number
2
\end{lstlisting}

\begin{keypoint}[Ask, then convert]
Ask with \code{s-exp-number?}, \code{s-exp-symbol?}, or \code{s-exp-list?}, then convert with the matching \code{s-exp->number}, \code{s-exp->symbol}, or \code{s-exp->list}. A conversion fails at run time if the guess was wrong.
\end{keypoint}

\begin{pitfall}[A lone quoted symbol is not an s-expression]
\code{'x} has type \tyname{symbol}, not \tyname{s-expression}, so \code{(parse 'x)} is a type error. To get an s-expression for one symbol, write \rkt|`x| (backquote) or \rkt|(symbol->s-exp 'x)|.
\end{pitfall}

\subsection{Concrete and Abstract Syntax}

The first interpreter is a calculator:
\begin{lstlisting}[style=repl]
'(+ 2 3)                     ===> 5
'(* 3 4)                     ===> 12
'(* (+ 2 3) (* 5 (+ 6 7)))   ===> 325
\end{lstlisting}
Following \cite[Chapters~2--3]{PLAI}, the work is split into separate steps (\Cref{fig:pipeline}):
\begin{itemize}
  \item \textbf{Concrete syntax} is what the programmer types, the text \code{(+ 2 3)}. Racket's \emph{reader} already turns text into an s-expression, so we get that step for free; quote is how we hand an s-expression over directly.
  \item \textbf{Abstract syntax} is a tree that keeps only the meaning. \code{parse} turns an s-expression into this tree, rejecting anything that is not a valid program.
  \item The \textbf{interpreter} walks the tree and computes the answer. \cite{PLAI} calls it \code{interp}; in lecture it is \code{eval}.
\end{itemize}

\begin{figure}[htbp]
\centering
\begin{tikzpicture}[node distance=4mm and 14mm]
  \node[box, fill=greybg] (t) {text};
  \node[box, fill=uspace, right=of t] (s) {s-expression};
  \node[box, fill=kspace, right=of s] (a) {\tyname{ArithE}};
  \node[box, fill=tsk, right=of a] (n) {number};
  \draw[arr] (t) -- node[above, lbl] {read} (s);
  \draw[arr] (s) -- node[above, lbl] {\code{parse}} (a);
  \draw[arr] (a) -- node[above, lbl] {\code{eval}} (n);
  \node[below=2mm of t, lbl] (et) {\code{"(+ 2 3)"}};
  \node[below=2mm of s, lbl] (es) {\code{'(+ 2 3)}};
  \node[below=2mm of a, lbl, align=center] (ea) {\code{(plusC (numC 2)}\\\code{(numC 3))}};
  \node[below=2mm of n, lbl] (en) {\code{5}};
  \draw[decorate, decoration={brace, amplitude=4pt, mirror}] ([yshift=-11mm]t.south west) -- ([yshift=-11mm]s.south east) node[midway, below=5pt, lbl] {concrete syntax};
  \draw[decorate, decoration={brace, amplitude=4pt, mirror}] ([yshift=-11mm,xshift=-6mm]a.south west) -- ([yshift=-11mm,xshift=6mm]a.south east) node[midway, below=5pt, lbl] {abstract syntax};
\end{tikzpicture}
\caption{The interpreter pipeline. Racket's reader provides the first step.}
\label{fig:pipeline}
\end{figure}

\subsection{The Calculator}

\subsubsection{Abstract syntax} One variant per kind of expression:
\begin{lstlisting}
(define-type ArithE
  [numC (n : number)]
  [plusC (e1 : ArithE) (e2 : ArithE)]
  [timesC (e1 : ArithE) (e2 : ArithE)])
\end{lstlisting}
The operands of \code{plusC} are \tyname{ArithE}, not \code{number}. That is what allows \code{'(* (+ 2 3) 4)}: the left operand of the multiplication is itself an addition, the same trick as \tyname{Tree} having subtrees (\Cref{fig:calcast}).

\begin{figure}[htbp]
\centering
\begin{tikzpicture}[level distance=9mm, every node/.style={ast},
    level 1/.style={sibling distance=34mm}, level 2/.style={sibling distance=17mm}, level 3/.style={sibling distance=12mm}]
  \node[asth] {timesC}
    child {node {plusC} child {node[astl] {numC 2}} child {node[astl] {numC 3}}}
    child {node[asth] {timesC}
      child {node[astl] {numC 5}}
      child {node {plusC} child {node[astl] {numC 6}} child {node[astl] {numC 7}}}};
\end{tikzpicture}
\caption{The abstract syntax of \code{'(* (+ 2 3) (* 5 (+ 6 7)))}. Parentheses are gone; the nesting is the tree.}
\label{fig:calcast}
\end{figure}

\subsubsection{The interpreter} One clause per variant:
\begin{lstlisting}
(define (eval (e : ArithE)) : number
  (type-case ArithE e
    [numC (n) n]
    [plusC (e1 e2) (+ (eval e1) (eval e2))]
    [timesC (e1 e2) (* (eval e1) (eval e2))]))
\end{lstlisting}
Three variants, three clauses, and a recursive call wherever the type mentions \tyname{ArithE}. \code{eval} never checks whether its input makes sense, and it does not have to: a value of type \tyname{ArithE} can only be built from the three constructors, so every tree it receives is well formed.

\subsubsection{The parser} That guarantee has to come from somewhere, and it comes from \code{parse}. The version written in lecture:
\begin{lstlisting}
(define (parse (s : s-expression)) : ArithE
  (cond
    [(s-exp-number? s) (numC (s-exp->number s))]
    [(s-exp-list? s)
     (let [(l (s-exp->list s))]
       (case (s-exp->symbol (first l))
         [(+) (plusC (parse (second l)) (parse (third l)))]
         [(*) (timesC (parse (second l)) (parse (third l)))]))]))
\end{lstlisting}
A number becomes a \code{numC}. A list is dispatched on its first element with \code{case}, which compares a symbol against literal symbols, and each operand is parsed recursively. Putting the steps together, \rkt|(eval (parse '(* (+ 2 3) (* 5 (+ 6 7)))))| is 325.

\subsection{What Malformed Input Does}

The lecture parser is right on well-formed input, but on bad input it either crashes with a message about something else, or quietly builds a wrong tree:
\begin{center}
\begin{tabular}{@{}ll@{}}
\toprule
input & result \\
\midrule
\code{(parse '(- 1 2))} & \code{case: no matching clause} \\
\code{(parse '(+ 1))} & \code{third: list contains too few elements} \\
\code{(parse '(+ 1 2 3))} & \code{(plusC (numC 1) (numC 2))}: the 3 is silently dropped \\
\bottomrule
\end{tabular}
\end{center}
The last one is the worst kind of bug: no error, just a program that means something other than what was written. A more careful parser checks the shape of each form and reports what was expected. This version is not from lecture, but every check in it addresses one row of the table:
\begin{lstlisting}
(define (parse (s : s-expression)) : ArithE
  (cond
    [(s-exp-number? s) (numC (s-exp->number s))]
    [(s-exp-list? s)
     (let [(l (s-exp->list s))]
       (if (and (= (length l) 3) (s-exp-symbol? (first l)))
           (case (s-exp->symbol (first l))
             [(+) (plusC (parse (second l)) (parse (third l)))]
             [(*) (timesC (parse (second l)) (parse (third l)))]
             [else (error 'parse "unknown operator")])
           (error 'parse "expected (op e1 e2)")))]
    [else (error 'parse "expected a number or a list")]))
\end{lstlisting}
\begin{itemize}
  \item The length check rejects \code{'(+ 1)} and \code{'(+ 1 2 3)}.
  \item \code{s-exp-symbol?} rejects \code{'(1 2 3)}, whose head is not an operator at all.
  \item The \code{else} of \code{case} rejects operators the language does not have, such as \code{-}.
  \item The final \code{else} rejects anything that is neither a number nor a list, such as a lone symbol.
\end{itemize}

\subsection{(SUPPLEMENT) Further Topics}

\subsubsection{Why parsing is cheap here} Real languages need a lexer and a parser, often generated from a grammar, to get from text to a tree. Lisp-family languages put that work into one generic step: the reader turns any text into nested lists of atoms, and each language then only has to interpret those lists. \cite[Chapter~2]{PLAI} uses this to skip parsing entirely. The flip side is that the s-expression \code{'(+ 1 2 3)} is perfectly well formed as data; deciding that it is not a valid \emph{calculator} program is the job of \code{parse}.

\subsubsection{Homoiconicity and quasiquote} Racket code is itself written as s-expressions, so programs can build programs as data. Quasiquote makes that convenient: if \code{x} is 5, then \rkt|`(+ ,x 1)| is the list \code{'(+ 5 1)}. Everything is quoted except the parts marked with a comma, which are evaluated and spliced in; \code{,@xs} splices in the elements of a whole list. Racket macros are functions from syntax to syntax built on the same idea, which is a large part of why Racket is good at implementing languages.

\subsubsection{Desugaring} The calculator has no subtraction. One way to add it without touching \code{eval} is to translate it in the parser: \code{(- a b)} becomes \code{(plusC a (timesC (numC -1) b))}. Languages are usually defined as a small \emph{core} interpreted directly plus a \emph{surface} that desugars into it~\cite[Chapter~4]{PLAI}. The core stays small enough to reason about, and the surface stays convenient.

\section{Recitation 3: Typed Boolean Expressions and Pattern Matching}
\label{sec:rec3}

Recitation~3 has two halves. Part~A (handout 3a) redoes Recitation~2's interpreter and De Morgan optimizer in \code{plai-typed}, to see what the type checker buys. Part~B (handout 3b) returns to plain Racket and studies \code{match}: literal patterns, wildcards, and taking lists apart.

\subsection{Part A: Boolean Expressions, Typed}

\subsubsection{From four structs to one type} In \Cref{sec:rec2}, a Boolean expression was built from four unrelated structs; ``a bexp'' existed only in our heads, and no field had a type. With \code{define-type}, the four become variants of \emph{one} type, and every field says what it holds:
\begin{lstlisting}
(define-type bexp-b
  [leaf-node-b (b : boolean)]
  [and-node-b (b1 : bexp-b) (b2 : bexp-b)]
  [or-node-b (b1 : bexp-b) (b2 : bexp-b)]
  [not-node-b (b : bexp-b)])
\end{lstlisting}

\begin{exercise}[Typed evaluator]
Implement \code{evaluate-bexp-b} with \code{type-case}.
\end{exercise}
\begin{solution}
\code{type-case} plays the role \code{match} played last week:
\begin{lstlisting}
(define (evaluate-bexp-b [b : bexp-b]) : boolean
  (type-case bexp-b b
    [leaf-node-b (v) v]
    [and-node-b (b1 b2) (and (evaluate-bexp-b b1) (evaluate-bexp-b b2))]
    [or-node-b (b1 b2) (or (evaluate-bexp-b b1) (evaluate-bexp-b b2))]
    [not-node-b (b1) (not (evaluate-bexp-b b1))]))

(test (evaluate-bexp-b (leaf-node-b #t)) #t)
(test (evaluate-bexp-b (and-node-b (leaf-node-b #t) (leaf-node-b #f))) #f)
(test (evaluate-bexp-b (not-node-b (or-node-b (leaf-node-b #f) (leaf-node-b #f)))) #t)
\end{lstlisting}
\code{test} is \code{plai-typed}'s \code{check-expect}, and \rkt|(print-only-errors true)| at the top of the file makes passing tests silent. This is the same shape as \code{eval} for \tyname{ArithE} in \Cref{sec:calc}.
\end{solution}

\subsubsection{What do the types buy us?}
\begin{enumerate}[label=(\alph*)]
  \item \textbf{Bad leaves are rejected before anything runs.} Last week's \code{bexp-sym-1} was
\begin{lstlisting}
(and-node (leaf-node 'x) (leaf-node 'y))
\end{lstlisting}
  and \code{evaluate-bexp} happily returned \code{'y}: Racket treats every value except \code{\#f} as true, so \rkt|(and 'x 'y)| is \code{'y}. No error, just an answer of the wrong kind. The typed version
\begin{lstlisting}
(evaluate-bexp-b (and-node-b (leaf-node-b 'x) (leaf-node-b 'y)))
\end{lstlisting}
  is rejected (\code{typecheck failed: boolean vs.\ symbol}) before anything runs, and the error points at the leaf that is wrong, not at some later function that trips over it.
  \item \textbf{Missing cases are rejected too.} Last week's \code{evaluate-bexp-m} handled only leaves; it was accepted and failed only when called on something else. The typed version
\begin{lstlisting}
(define (evaluate-bexp-b-m [b : bexp-b]) : boolean
  (type-case bexp-b b
    [leaf-node-b (v) v]))
\end{lstlisting}
  stops the program from starting: \code{you did not include a case for the and-node-b variant}. This is the \code{area-v3}/\code{area-v4} contrast of \Cref{sec:definetype} again.
\end{enumerate}

\subsubsection{De Morgan, typed} Now the leaves are variables:
\begin{lstlisting}
(define-type bexp
  [leaf-node (s : symbol)]
  [and-node (b1 : bexp) (b2 : bexp)]
  [or-node (b1 : bexp) (b2 : bexp)]
  [not-node (b1 : bexp)])
\end{lstlisting}

\begin{exercise}[Port \code{demorgan-opt} to \code{plai-typed}]
Rewrite the optimizer of \Cref{sec:rec2} for the type \tyname{bexp}.
\end{exercise}
\begin{solution}
Last week's \code{match} could name a nested pattern directly, such as \rkt|(not-node (and-node b1 b2))|. A \code{type-case} clause only takes apart the \emph{outermost} constructor, so the \code{not-node} clause needs a second \code{type-case}, on the thing being negated:
\begin{lstlisting}
(define (demorgan-opt [b : bexp]) : bexp
  (type-case bexp b
    [leaf-node (v) (leaf-node v)]
    [and-node (b1 b2) (and-node (demorgan-opt b1) (demorgan-opt b2))]
    [or-node (b1 b2) (or-node (demorgan-opt b1) (demorgan-opt b2))]
    [not-node (b1)
      (type-case bexp b1
        [leaf-node (v) (not-node (leaf-node v))]
        [and-node (b2 b3) (or-node (demorgan-opt (not-node b2))
                                   (demorgan-opt (not-node b3)))]
        [or-node (b2 b3) (and-node (demorgan-opt (not-node b2))
                                   (demorgan-opt (not-node b3)))]
        [not-node (b2) (demorgan-opt b2)])]))
\end{lstlisting}
The inner clauses are the four nested patterns of last week, one per variant of the negated expression: a negated leaf stays; $\neg(b_2 \land b_3)$ becomes $\neg b_2 \lor \neg b_3$ with both new negations optimized again; $\neg(b_2 \lor b_3)$ symmetrically; and $\neg\neg b_2$ becomes the optimized $b_2$. Both \code{type-case}s are checked for exhaustiveness, so a forgotten shape of negation is now a static error.
\end{solution}

The tests no longer need \code{syn-equal?}: \code{test} compares \code{define-type} values structurally.
\begin{lstlisting}
(define bexp-1 (not-node (or-node (not-node (leaf-node 'P)) (leaf-node 'Q))))
(define bexp-2 (or-node (not-node (not-node (leaf-node 'P))) (not-node (leaf-node 'Q))))
(define bexp-3 (or-node (leaf-node 'P) (not-node (leaf-node 'Q))))

(test (demorgan-opt bexp-1) (and-node (leaf-node 'P) (not-node (leaf-node 'Q))))
(test (demorgan-opt bexp-2) (or-node (leaf-node 'P) (not-node (leaf-node 'Q))))
(test (demorgan-opt bexp-3) (or-node (leaf-node 'P) (not-node (leaf-node 'Q))))
\end{lstlisting}

\subsection{Part B: Pattern Matching with \texttt{match}}

So far \code{match} has only taken structs apart. It can also match literal values, ignore parts of the input with wildcards, and take lists apart.

\subsubsection{Matching on values} A pattern can be a literal, which makes \code{match} a compact conditional:
\begin{lstlisting}
(define (factorial x)
  (match x
    [0 1]
    [_ (* x (factorial (- x 1)))]))

(define (factorial-if x)
  (if (= x 0)
      1
      (* x (factorial-if (- x 1)))))
\end{lstlisting}
\code{\_} is the \emph{wildcard}: it matches anything and binds nothing. Clauses are tried top to bottom, so \code{\_} goes last as the catch-all. With two cases \code{match} and \code{if} are about the same size; with more, nested \code{if}s pile up:
\begin{lstlisting}
(define (fibonacci-if x)
  (if (= x 0)
      0
      (if (= x 1)
          1
          (+ (fibonacci-if (- x 1)) (fibonacci-if (- x 2))))))
\end{lstlisting}

\begin{exercise}[Fibonacci with \code{match}]
Rewrite \code{fibonacci-if} using \code{match}.
\end{exercise}
\begin{solution}
One line per case, and the structure stays flat:
\begin{lstlisting}
(define (fibonacci x)
  (match x
    [0 0]
    [1 1]
    [_ (+ (fibonacci (- x 1)) (fibonacci (- x 2)))]))

(check-expect (fibonacci 1) 1)
(check-expect (fibonacci 10) 55)
(check-expect (fibonacci 10) (fibonacci-if 10))
\end{lstlisting}
\end{solution}

\begin{exercise}[What is \code{else}?]
\code{[else ...]} often appears as the last clause of a \code{match}. Predict \rkt|(match 5 [else else])|.
\end{exercise}
\begin{solution}
It is \code{5}. Inside \code{match}, \code{else} is not a keyword but a \emph{pattern variable}, like \code{x}: it matches anything and is bound to what it matched. So here \code{else} is bound to 5 and the clause returns it. That is why \code{[else ...]} happens to work as a catch-all; but it is an accident, and it quietly binds a variable named \code{else} in the body. Use \code{\_} instead.
\end{solution}

\subsubsection{Destructuring lists} Two attempts at ``the third element of a list'':
\begin{lstlisting}
(define (my-third x)
  (match x
    [(list a b c d) c]
    [_ x]))

(define (other-third lst)
  (match lst
    [(cons _ (cons _ (cons third _))) third]
    [_ (error "List too short")]))
\end{lstlisting}

\begin{exercise}[Predict]
On which inputs do the two functions behave differently, and why?
\end{exercise}
\begin{solution}
\rkt|(list a b c d)| matches a list of \emph{exactly} four elements. \rkt|(cons _ (cons _ (cons third _)))| matches any list of \emph{at least} three: the last \code{\_} matches whatever rest is left, including \code{'()}.
\begin{center}
\begin{tabular}{@{}lll@{}}
\toprule
input & \code{my-third} & \code{other-third} \\
\midrule
\rkt|(list 2 3)| & \code{'(2 3)} & error \code{"List too short"} \\
\code{0} & \code{0} & error \\
\rkt|(list '(2 3) 2 4 5 6 7)| & the whole list & \code{4} \\
\rkt|(list 1 2 3 4)| & \code{3} & \code{3} \\
\bottomrule
\end{tabular}
\end{center}
\begin{lstlisting}
(check-expect (my-third (list 2 3)) (list 2 3))
(check-expect (my-third 0) 0)
(check-expect (my-third (list '(2 3) 2 4 5 6 7)) (list '(2 3) 2 4 5 6 7))
(check-expect (my-third (list 1 2 3 4)) 3)

(check-error (other-third (list 2 3)) "List too short")
(check-error (other-third 0) "List too short")
(check-expect (other-third (list '(2 3) 2 4 5 6 7)) 4)
(check-expect (other-third (list 1 2 3 4)) 3)
\end{lstlisting}
\code{my-third} is right only on four-element lists, and when it is wrong it does not fail: the fallback \rkt|[_ x]| hands back its input, so the caller gets a list or a number where it expected the third element. That is the same kind of silent wrong answer as the calculator parser turning \code{'(+ 1 2 3)} into $1 + 2$ (\Cref{sec:calc}). Raising an error, as \code{other-third} does, is better.
\end{solution}

\begin{exercise}[Fix \code{my-third}]
Make \code{my-third} return the third element of any list with at least three elements and raise an error otherwise, using a \code{list} pattern. Hint: \code{\_ ...} matches zero or more elements.
\end{exercise}
\begin{solution}
\begin{lstlisting}
(define (my-third-fixed x)
  (match x
    [(list _ _ c _ ...) c]
    [_ (error "List too short")]))

(check-error (my-third-fixed (list 2 3)) "List too short")
(check-expect (my-third-fixed (list 1 2 3)) 3)
(check-expect (my-third-fixed (list '(2 3) 2 4 5 6 7)) 4)
\end{lstlisting}
This is equivalent to \code{other-third}'s nested \code{cons} pattern, and easier to read.
\end{solution}

\subsection{\texttt{match} vs.\ \texttt{type-case}}

The two halves used two tools for taking values apart. Each is strong where the other is weak:
\begin{center}
\begin{tabular}{@{}lll@{}}
\toprule
 & \code{match} & \code{type-case} \\
\midrule
language & \code{racket} & \code{plai-typed} \\
nested patterns & yes: \rkt|(not-node (and-node b1 b2))| & no: one level; nest another \code{type-case} \\
literals, wildcards, lists & yes & no \\
missing case & run-time error, only if reached & rejected before the program runs \\
\bottomrule
\end{tabular}
\end{center}
\code{match} is more expressive. \code{type-case} is more limited, and in exchange it can check, before anything runs, that every variant is handled. Recitation~4 (\Cref{sec:rec4}) implements a small \code{match} in our own interpreter.

\subsection{(SUPPLEMENT) Further Topics}

\subsubsection{More \code{match} patterns} Racket's \code{match} has many more pattern forms~\cite[Section~12]{RacketGuide}: \rkt|(? even? n)| matches when the predicate holds and binds \code{n}; \rkt|(and p1 p2)| and \rkt|(or p1 p2)| combine patterns; \rkt|(list-rest a b rest)| names the tail; \rkt|(app f p)| matches \code{p} against \rkt|(f v)|; and a clause can carry a guard with \code{\#:when}. Repeating a pattern variable, as in \rkt|(list x x)|, requires the two parts to be \code{equal?}.

\subsubsection{Exhaustiveness for richer patterns} Languages with both nested patterns and types, such as OCaml, Haskell, and Rust, get the best of both columns of the table: the compiler checks a set of nested patterns for exhaustiveness and also warns about \emph{redundant} clauses that can never match (like a clause after \code{\_}). The check is done by computing which values are not yet covered, clause by clause, and it is the same algorithm a compiler uses to turn a \code{match} into an efficient decision tree.

